% Part: sets-functions-relations
% Chapter: size-of-sets
% Section: introduction
% Hindi translation (hi-Deva-IN) of Open Logic commit 9620cc73f9c8e0ad003c514a5d3748f29611c4c0.

\documentclass[../../../include/open-logic-section]{subfiles}

\begin{document}

\olfileid[hi]{sfr}{siz}{int}

\olsection{परिचय}

जब जॉर्ज कैंटर ने 1870 के दशक में समुच्चय सिद्धांत विकसित किया, तो उनका एक
उद्देश्य अनंत संग्रह---मध्ययुगीन विद्वानों की शब्दावली में वास्तविक
अनंतता---की धारणा को स्वीकार्य बनाना था। इसका एक महत्त्वपूर्ण अंग विभिन्न
समुच्चयों के \emph{आकार} का उनका निरूपण था। यदि $a$, $b$ और $c$ सभी भिन्न
हैं, तो सहज बोध से समुच्चय $\{a, b, c\}$, $\{a, b\}$ से \emph{बड़ा} है।
पर अनंत समुच्चयों का क्या? क्या वे सभी एक-दूसरे जितने बड़े हैं? पता चलता है
कि ऐसा नहीं है।

यहाँ पहली महत्त्वपूर्ण अवधारणा गणन की है। हम प्रत्येक परिमित समुच्चय के सभी
!!{element}s लिखकर उसकी सूची बना सकते हैं। कुछ अनंत समुच्चयों के भी सभी
!!{element}s सूचीबद्ध किए जा सकते हैं, यदि सूची को स्वयं अनंत होने दें।
ऐसे समुच्चय !!{enumerable} कहलाते हैं। कैंटर का आश्चर्यजनक परिणाम---जिसे हम
इस अध्याय के अंत तक पूरी तरह समझ लेंगे---यह था कि कुछ अनंत समुच्चय
!!{enumerable} नहीं होते।

\end{document}
